\uuid{PJHa}
\titre{Solution générale/solution particulière/condition initiale}
\auteur{hanani ; mzari}
\organisation{exo7}
\licence{CC BY-NC-SA 4.0}
\datecreate{2020-01-20}
\niveau{L1}
\module{Analyse}
\chapitre{Equation différentielle}
\sousChapitre{Résolution d'équation différentielle du premier ordre}
\difficulte{3}
\theme{Equations du premier ordre, solution générale, solution particulière, condition initiale}
\exo7qcmid{lille-2.equations-differentielles.008}

\contenu{
\begin{question}
\qclassification{225.01}
\qtags{motcle=solution générale/solution particulière/condition initiale}
On considère l'équation différentielle $(E)$ : $\displaystyle (1+x^2)y'-y=1$. Parmi les affirmations suivantes cocher celles qui sont vraies :
\begin{answers}
\bad{La solution générale de l'équation homogène est : $\displaystyle y=k\mathrm{e}^{-\arctan x}$, $k\in \Rr$.}
\good{La solution générale de $(E)$ est : $\displaystyle y=-1+k\mathrm{e}^{\arctan x}$, $k\in \Rr$.}
\good{L'unique solution de $(E)$ vérifiant $y(0)=0$ et $\displaystyle y=-1+\mathrm{e}^{\arctan x}$.}
\bad{$(E)$ n'admet pas de solution vérifiant $y'(0)=0$.}
\end{answers}
\begin{explanations}
Les solutions sur $\Rr$ de l'équation homogène sont les fonctions $\displaystyle Y=k\mathrm{e}^{\arctan x}$, $k\in \Rr$, et $y_0=-1$ est une solution particulière. Les solutions de $(E)$ sont donc les fonctions
$$y=-1+k\mathrm{e}^{\arctan x},\quad k\in \Rr.$$
La condition $y(0)=0$ donne $k=1$. Par ailleurs, $y=-1$ est une solution de $(E)$ vérifiant $y'(0)=0$.
\end{explanations}
\end{question}
}
